\section{Related Work}
\label{sec:related}

\paragraph{Aggregate uncertainty from unresolved relationships.}
Resolution-aware query answering retains unresolved entities and computes
aggregate bounds without requiring probabilities over resolutions
\citep{sismanis2009resolution}. Probabilistic linkage supports aggregate
queries \citep{hua2012aggregate}; constrained probabilistic similarity joins
likewise define one-to-one, one-to-many, and many-to-many matching worlds and
evaluate aggregates over them \citep{chen2019constrained}. Bayesian linkage
represents entire compatible matchings and quantifies their uncertainty
\citep{sadinle2017linkage}. Our closest optimization comparator,
\citet{turkcapar2023aggregate}, computes aggregate ranges over
candidate-constrained one-to-many table integrations via graph matching. It
also studies false-negative candidates and failure to cover the true
aggregate. Thus possible-world aggregation, matching constraints, extremal
bounds, and candidate-omission risk are not new here.

The narrow distinction is the latent world being quantified over. A matching
selects links; our world partitions rows into variable-cardinality temporal
events. Each event must admit one common timestamp completion, connected
positive overlap, and bounded occupancy at every instant. Consequently,
membership can exceed capacity through sequential turnover, while individually
compatible links may be jointly unrealizable. These are object-level
constraints rather than a new semantics for uncertain aggregation.

\paragraph{Possible worlds, repairs, and sensitivity.}
Range-consistent query answering optimizes aggregates across legal database
repairs, with tractability results for query rewriting under primary-key
violations \citep{elkhalfioui2024range,elkhalfioui2026upper}.
Attribute-uncertain databases propagate value bounds through relational
queries \citep{feng2021audb}. EventFrontier adopts the same general principle
of quantifying over admissible worlds; its contribution concerns a temporal
event family, rather than a new semantics for uncertain data. Query resilience
studies the smallest input deletion that makes a Boolean query false
\citep{freire2015resilience}. Our omission budgets instead parameterize
candidate support for latent relationships. We report sensitivity within this
declared family, without claiming a new general notion of robustness or an
unconditional coverage guarantee.

\paragraph{Interval optimization and decomposition.}
Interval scheduling already admits structured pricing, column generation,
and branch-and-price. In particular, \citet{muir2023interval} partition fixed
interval jobs under economies-of-scale costs, giving polynomial pricing for
the max-weight case and pseudo-polynomial pricing for a broader class.
Their schedules forbid simultaneous jobs on the same machine; our events
require positive-overlap connectivity while limiting occupancy. Our local
result jointly encodes segment coverage, boundary bridges, capacity, and
fixed membership in a consecutive-ones formulation. For fixed exact times,
root, span, and additive weights, this yields an integral event-pricing LP.
Total unimodularity and decomposition are established tools; the specific
claim is their applicability to this connected event formulation.
The coupled master can remain fractional, and branch-compatible pricing may
require exponentially many cases. We make no polynomial-time claim for the
full partition problem or for arbitrary set-valued timestamps.

\paragraph{Mobility as an observation setting.}
Shareability networks and trip--vehicle assignment optimize prospective
allocations \citep{santi2014shareability,alonsomora2017ondemand}.
Studies of public ride-hailing data predict sharing outcomes or construct
proxy groups \citep{taiebat2022sharing,sebtichen2026hidden}.
Our application instead asks which retrospective aggregates are compatible
with released rows after event identifiers and time precision are removed.
The public audits establish conditional ranges, not recovery of actual
vehicle assignments or evidence of social interaction. Chicago and ATR
dyads test the matching special case and candidate-support validity; evidence
for the general temporal-event contribution must additionally exercise
sequential membership and joint timestamp feasibility.
